\subsection{Computing minimum palindromic length: Problem 16.9}
\label{cand:16.9}

In the free group \(F_n\) with specified basis \(x_1,\ldots,x_n\),
let \(\operatorname{pl}(w)\) be the minimum number of palindromes
whose product is \(w\), with \(\operatorname{pl}(1)=0\).
The following algorithm uses \(O(L^3)\) arithmetic operations
and \(O(L^2)\) integer storage for an input word of length \(L\).
Its stored integers have magnitude \(O(L)\).

Let \(\tau(x_i)=x_i^{-1}\), and form
\(H=F_n\rtimes\langle t\rangle\), with \(t^2=1\)
acting by \(\tau\). The substitutions \(s_0=t,s_i=x_it\)
give the presentation
\[
 H=\langle s_0,\ldots,s_n:s_i^2=1\text{ for all }i\rangle.
\]
The subgroup \(F_n\) is its even-parity subgroup.
An element \(p\in F_n\) is a palindrome exactly when
\(\tau(p)=p^{-1}\), hence when \(pt\) is an involution.
Every nonidentity involution of this free product
is conjugate to an \(s_i\), by cyclic reduction,
and has odd parity. Thus the involutions, or
reflections, are precisely the elements \(pt\)
for palindromes \(p\).

If \(r(h)\) is reflection length in \(H\), then
\begin{equation}\label{eq:palindrome-reflection}
 \operatorname{pl}(w)=\min\{r(w),r(wt)\}.
\end{equation}
Indeed, a product \(w=p_1\cdots p_k\) gives
\[
 wt^k=(p_1t)(\tau(p_2)t)\cdots
                   (\tau^{k-1}(p_k)t).
\]
Conversely, write a reflection factorization of
\(wt^j\), \(j=0,1\), as
\((q_1t)\cdots(q_mt)\). Parity gives \(m\equiv j\pmod2\),
so
\(w=q_1\tau(q_2)\cdots\tau^{m-1}(q_m)\);
each factor is a palindrome. This proves both
inequalities in \eqref{eq:palindrome-reflection}.

For a word \(V\) on the involution alphabet, let
\(d(V)\) be the minimum number of unmatched positions
in a noncrossing partial pairing of equal letters.
This is unchanged by insertion or deletion of two
adjacent equal letters. Pairing the inserted letters
gives one inequality. For the reverse, start with
a pairing before deletion. If the deleted positions
are paired together, remove that pair. If neither
has a partner, remove them. If just one has another
partner, leave that partner unmatched. If both have
other partners, pair those partners to each other.
Their labels agree, and noncrossing is preserved:
the old arcs joined along the short deleted interval
give a new arc disjoint from the others. None of
these repairs increases the number unmatched.

It follows from free-product reduction that \(d\)
depends only on the represented element. Moreover
\(d(h)=r(h)\). A conjugate \(vs_iv^{-1}\) has a
pairing leaving only its central letter unmatched,
so a product of \(k\) reflections has \(d\leq k\).
Conversely, delete the unmatched letters of a
pairing. The remaining fully paired word reduces
to the identity by successively removing innermost
pairs. Deleting one letter uses the identity
\[
 Xs_iY=(Xs_iX^{-1})(XY).
\]
Successive deletions express the original word
as a product of that many reflections. This
proves equality. The general Coxeter deletion
principle is prior work of Dyer~\cite{dyer};
the special case needed here has been proved directly.

Replace each \(x_i\) in the input by \(i0\) and each
\(x_i^{-1}\) by \(0i\), obtaining \(W\) of length
at most \(2L\). Then
\[
 \operatorname{pl}(w)=\min\{d(W),d(W0)\}.
\]
For half-open intervals of \(V=W0\), set \(D(i,i)=0\) and compute
\[
 D(i,j)=\min\left\{
 \begin{aligned}
 &1+D(i+1,j),\\
 &\min_{\substack{i<k<j\\V_i=V_k}}
       \bigl(D(i+1,k)+D(k+1,j)\bigr)
 \end{aligned}
 \right\}.
\]
The minimum over an empty set is \(+\infty\).
The first option leaves position \(i\) unmatched.
Pairing it with \(k\) separates the remaining
positions into the inside and the following
interval, by noncrossing. This exhausts all
possibilities, proving the recurrence. There
are \(O(L^2)\) intervals with \(O(L)\) choices each.

Stored minimizing choices recover the unmatched
positions, whose successive deletions recover
reflection factors \(r_i\). Put \(q_i=r_it\),
rewrite their even-parity words in the \(x_i\),
and apply \(\tau\) to every second factor as above.
This gives a minimum palindrome factorization.
The total output length is \(O(L^2)\), since
each conjugating prefix has length \(O(L)\)
and there are \(O(L)\) factors. Bit complexity
adds the usual \(O(\log(L+1))\) overhead for
the stored integers; encoding generator indices
is separate. This computes group palindromic
length, allowing cancellation between factors,
not merely a partition of the reduced word.
Source: \artifact{research/16.9-proof.md}.
